% GENERATED FILE. Edit canonical lessons, then run npm run generate:books.
% canonical-source-hash: fnv1a64:6fef5d496d60b3d4
% canonical-lessons: JA-C523-daredemo, JA-C523-doredemo, JA-C523-sazo, JA-C523-tamatama, JA-C523-nantoka

\chapter{Anyone, Any One}
\label{ch:ja-anyone-any-one}
\hlchaptermodality{523}

\begin{chapteropening}
\textbf{By the end of this chapter:} \emph{I can offer things and guess how they will turn out.}

Complete the last lesson of chapter 523: I can offer things and guess how they will turn out.
\end{chapteropening}

\section[\texorpdfstring{\ja{だれでも} (daredemo)}{だれでも (daredemo)}]{\ja{だれでも} (daredemo) — anyone; everyone}

\label{lesson:JA-C523-daredemo}



\begin{quote}
Before the new one: say the Japanese for at any time, then the Japanese for anything.
\end{quote}

\subsection*{You'll want to know: \ja{だれでも}}
\textbf{\ja{だれでも}} — \emph{daredemo} — \textquotedblleft{}anyone; everyone\textquotedblright{}.

\begin{quote}
\textbf{\ja{もし} \ja{じかんが} \ja{あったら、だれでも} \ja{きて} \ja{ください。}} — \emph{moshi jikan ga attara, daredemo kite kudasai} — \textquotedblleft{}If you have time, anyone is welcome to come.\textquotedblright{}
\end{quote}

\begin{cousinweb}
\textbf{\ja{だれでも}} (\emph{daredemo}) is a native Japanese word (wago), inherited rather than borrowed from Chinese: built from \emph{dare}, \textquotedblleft{}who\textquotedblright{}, and \emph{demo}, \textquotedblleft{}even\textquotedblright{} (JMdict).
\end{cousinweb}

\begin{grammarlens}[title={What you've built}]
One more word for saying what happens if.
\end{grammarlens}

\subsection*{Your turn}
\begin{itemize}
\raggedright
  \item \emph{Say it:} \emph{daredemo}
  \item \emph{Say it:} \emph{daredemo}, once more
  \item \emph{Say it:} \emph{moshi jikan ga attara, daredemo kite kudasai} — If you have time, anyone is welcome to come.
  \item \emph{Recall:} say the Japanese for at any time, then the Japanese for anything, then say \emph{daredemo} again
\end{itemize}

\begin{tcolorbox}[breakable,colback=teal!4,colframe=teal!35!black,title={Before you move on}]
What does \ja{だれでも} mean? (\textquotedblleft{}Anyone; everyone\textquotedblright{}.) Say it once more.
\end{tcolorbox}
\section[\texorpdfstring{\ja{どれでも} (doredemo)}{どれでも (doredemo)}]{\ja{どれでも} (doredemo) — any one (of them)}

\label{lesson:JA-C523-doredemo}



\begin{quote}
Before the new one: say the Japanese for anything, then the Japanese for anyone.
\end{quote}

\subsection*{You'll want to know: \ja{どれでも}}
\textbf{\ja{どれでも}} — \emph{doredemo} — \textquotedblleft{}any one (of them)\textquotedblright{}.

\begin{quote}
\textbf{\ja{もし} \ja{ほしい} \ja{おかしが} \ja{あったら、どれでも} \ja{たべて} \ja{ください。}} — \emph{moshi hoshii okashi ga attara, doredemo tabete kudasai} — \textquotedblleft{}If there are any sweets you want, eat whichever you like.\textquotedblright{}
\end{quote}

\begin{cousinweb}
\textbf{\ja{どれでも}} (\emph{doredemo}) is a native Japanese word (wago), inherited rather than borrowed from Chinese, and usually written in kana: built from \emph{dore}, \textquotedblleft{}which\textquotedblright{}, and \emph{demo}, \textquotedblleft{}even\textquotedblright{} (JMdict).
\end{cousinweb}

\begin{grammarlens}[title={What you've built}]
One more word for saying what happens if.
\end{grammarlens}

\subsection*{Your turn}
\begin{itemize}
\raggedright
  \item \emph{Say it:} \emph{doredemo}
  \item \emph{Say it:} \emph{doredemo}, once more
  \item \emph{Say it:} \emph{moshi hoshii okashi ga attara, doredemo tabete kudasai} — If there are any sweets you want, eat whichever you like.
  \item \emph{Recall:} say the Japanese for anything, then the Japanese for anyone, then say \emph{doredemo} again
\end{itemize}

\begin{tcolorbox}[breakable,colback=teal!4,colframe=teal!35!black,title={Before you move on}]
What does \ja{どれでも} mean? (\textquotedblleft{}Any one (of them)\textquotedblright{}.) Say it once more.
\end{tcolorbox}
\section[\texorpdfstring{\ja{さぞ} (sazo)}{さぞ (sazo)}]{\ja{さぞ} (sazo) — no doubt; I imagine}

\label{lesson:JA-C523-sazo}



\begin{quote}
Before the new one: say the Japanese for anyone, then the Japanese for any one (of them).
\end{quote}

\subsection*{You'll want to know: \ja{さぞ}}
\textbf{\ja{さぞ}} — \emph{sazo} — \textquotedblleft{}no doubt; I imagine\textquotedblright{}.

\begin{quote}
\textbf{\ja{もし} \ja{かぞくで} \ja{りょこうしたら、さぞ} \ja{たのしいでしょう。}} — \emph{moshi kazoku de ryokō shitara, sazo tanoshii deshō} — \textquotedblleft{}If the family went on a trip together, it would surely be fun.\textquotedblright{}
\end{quote}

\begin{cousinweb}
\textbf{\ja{さぞ}} (\emph{sazo}) is usually counted a native Japanese word (wago), inherited rather than borrowed from Chinese: JMdict marks it as usually written in kana, with one rare kanji spelling, and no dictionary we use records an older form, so its history is not settled (JMdict).
\end{cousinweb}

\begin{grammarlens}[title={What you've built}]
One more word for saying what happens if.
\end{grammarlens}

\subsection*{Your turn}
\begin{itemize}
\raggedright
  \item \emph{Say it:} \emph{sazo}
  \item \emph{Say it:} \emph{sazo}, once more
  \item \emph{Say it:} \emph{moshi kazoku de ryokō shitara, sazo tanoshii deshō} — If the family went on a trip together, it would surely be fun.
  \item \emph{Recall:} say the Japanese for anyone, then the Japanese for any one (of them), then say \emph{sazo} again
\end{itemize}

\begin{tcolorbox}[breakable,colback=teal!4,colframe=teal!35!black,title={Before you move on}]
What does \ja{さぞ} mean? (\textquotedblleft{}No doubt; I imagine\textquotedblright{}.) Say it once more.
\end{tcolorbox}
\section[\texorpdfstring{\ja{たまたま} (tamatama)}{たまたま (tamatama)}]{\ja{たまたま} (tamatama) — as it happens; by accident}

\label{lesson:JA-C523-tamatama}



\begin{quote}
Before the new one: say the Japanese for any one (of them), then the Japanese for no doubt.
\end{quote}

\subsection*{You'll want to know: \ja{たまたま}}
\textbf{\ja{たまたま}} — \emph{tamatama} — \textquotedblleft{}as it happens; by accident\textquotedblright{}.

\begin{quote}
\textbf{\ja{もし} \ja{たまたま} \ja{じかんが} \ja{あったら、あそびに} \ja{きて} \ja{ください。}} — \emph{moshi tamatama jikan ga attara, asobi ni kite kudasai} — \textquotedblleft{}If you happen to have time, come and see us.\textquotedblright{}
\end{quote}

\begin{cousinweb}
\textbf{\ja{たまたま}} (\emph{tamatama}) is a native Japanese word (wago), inherited rather than borrowed from Chinese, and usually written in kana: \emph{tama}, the native reading KANJIDIC2 gives the kanji for \textquotedblleft{}accidentally\textquotedblright{}, said twice (JMdict, KANJIDIC2).
\end{cousinweb}

\begin{grammarlens}[title={What you've built}]
One more word for saying what happens if.
\end{grammarlens}

\subsection*{Your turn}
\begin{itemize}
\raggedright
  \item \emph{Say it:} \emph{tamatama}
  \item \emph{Say it:} \emph{tamatama}, once more
  \item \emph{Say it:} \emph{moshi tamatama jikan ga attara, asobi ni kite kudasai} — If you happen to have time, come and see us.
  \item \emph{Recall:} say the Japanese for any one (of them), then the Japanese for no doubt, then say \emph{tamatama} again
\end{itemize}

\begin{tcolorbox}[breakable,colback=teal!4,colframe=teal!35!black,title={Before you move on}]
What does \ja{たまたま} mean? (\textquotedblleft{}As it happens; by accident\textquotedblright{}.) Say it once more.
\end{tcolorbox}
\section[\texorpdfstring{\ja{なんとか} (nantoka)}{なんとか (nantoka)}]{\ja{なんとか} (nantoka) — somehow (or other)}

\label{lesson:JA-C523-nantoka}



\begin{quote}
Before the new one: say the Japanese for no doubt, then the Japanese for as it happens.
\end{quote}

\subsection*{You'll want to know: \ja{なんとか}}
\textbf{\ja{なんとか}} — \emph{nantoka} — \textquotedblleft{}somehow (or other)\textquotedblright{}.

\begin{quote}
\textbf{\ja{もし} \ja{がんばったら、かれは} \ja{なんとか} \ja{まにあうでしょう。}} — \emph{moshi ganbattara, kare wa nantoka maniau deshō} — \textquotedblleft{}If he tries hard, he will make it somehow.\textquotedblright{}
\end{quote}

\begin{cousinweb}
\textbf{\ja{なんとか}} (\emph{nantoka}) is a native Japanese word (wago), inherited rather than borrowed from Chinese, and usually written in kana: JMdict spells it with the kanji of \emph{nan}, \textquotedblleft{}what\textquotedblright{}, followed by the particles \emph{to} and \emph{ka} (JMdict).
\end{cousinweb}

\begin{grammarlens}[title={What you've built}]
One more word for saying what happens if.
\end{grammarlens}

\subsection*{Your turn}
\begin{itemize}
\raggedright
  \item \emph{Say it:} \emph{nantoka}
  \item \emph{Say it:} \emph{nantoka}, once more
  \item \emph{Say it:} \emph{moshi ganbattara, kare wa nantoka maniau deshō} — If he tries hard, he will make it somehow.
  \item \emph{Recall:} say the Japanese for no doubt, then the Japanese for as it happens, then say \emph{nantoka} again
\end{itemize}

\begin{tcolorbox}[breakable,colback=teal!4,colframe=teal!35!black,title={Before you move on}]
What does \ja{なんとか} mean? (\textquotedblleft{}Somehow (or other)\textquotedblright{}.) Say it once more.
\end{tcolorbox}
